Use \wikilink{s07--ed-centers--class-equation}{the class equation} and
\wikilink{s08--b007}{Sylow existence}. The second proof below uses Sylow
existence, so it must not be used as a premise in proving that theorem.

\begin{theorem}\label{p-group-center}
Every nontrivial finite $p$-group has nontrivial center; indeed
$p\mid|Z(G)|$.
\end{theorem}
\begin{proof}\label{p-group-center-proof}
The index of a proper subgroup of a finite $p$-group is a positive
power of $p$. In the class equation every noncentral centralizer is
proper, so all noncentral class sizes are divisible by $p$.
Since $p\mid|G|$, the equation gives $p\mid|Z(G)|$. The center contains
the identity, so its size is at least $p$.
\end{proof}
\begin{theorem}\label{cauchy-theorem}
If $G$ is finite and a prime $p$ divides $|G|$, then $G$ has an element
of order $p$.
\end{theorem}
\begin{proof}\label{cauchy-theorem-proof}
A Sylow $p$-subgroup $P$ is nontrivial. Choose $x\in P$ different
from the identity. By Lagrange its order is $p^r$ for some $r\ge1$.
Then $x^{p^{r-1}}$ has order $p$: its $p$th power is the identity,
but it is not itself the identity by minimality of the order of $x$.
\end{proof}
\begin{remark}\label{p-center-scope}
The center assertion excludes the trivial group. The convention here
is the finite prime-power-order definition in
\wikilink{s08--b001}{the source definition of a p-group}.
\end{remark}
